\subsection{Effective cyclic torsion and centralizer recognition in free products}

\begin{result}{A.4.7}{Proved}{One homology cover recognizes the cyclic data at a prescribed finite level}
Let
\[
 E=A_1*\cdots*A_k*F_r
\]
be the specified free product of finite groups and a finite-rank free
group.  Let $q:E\twoheadrightarrow Q$ be a specified finite quotient
injective on each finite factor, and put $K=\ker q$.
For a prime $p$, define the actual finite quotient
\[
 R_p=E/(K'K^p),\qquad
 |R_p|=|Q|p^d,\qquad
 d=1-|Q|\left(\sum_{i=1}^k\frac1{|A_i|}-k-r+1\right).
                                                               \tag{FG.1}
\]
For an element $a\in A_i$ of order $p$, the images in $Q$ of its
centralizer and cyclic-subgroup normalizer in $R_p$ are exactly
\[
 q(C_{A_i}(a)),\qquad q(N_{A_i}(\langle a\rangle)),
                                                               \tag{FG.2}
\]
respectively.  Every order-$p$ element of $R_p$ with nonidentity
image in $Q$ has image conjugate in $Q$ to the image of an
order-$p$ element of one of the $A_i$.
These statements describe images in the prescribed quotient $Q$;
they do not assert that all order-$p$ elements in $R_p$ are
conjugate to actual vertex-factor elements in $R_p$.
\end{result}

\begin{proof}
Use the Bass--Serre tree $T$ of the displayed free product, with
trivial edge stabilizers.  Faithfulness on the finite factors makes
$K$ act freely.  Thus $K$ is the fundamental group of the finite
graph $Y=K\backslash T$.  The Euler-characteristic calculation gives
its free rank $d$ in (FG.1), and its elementary abelian quotient
$K/K'K^p$ is literally $H_1(Y,\mathbb F_p)$.  This proves the order
formula and identifies its homology kernel.

Choose a vertex $v$ of $T$ with stabilizer $A_i$ and the corresponding
vertex of $Y$.  Write $\bar a=q(a)$.  The nontrivial group
$\langle\bar a\rangle$ acts freely on the oriented edges of $Y$,
because edges have trivial stabilizer in its deck group $Q$.
It preserves edge orientations.  Consequently an invariant
$\mathbb F_p$ edge-chain is a sum of complete $p$-edge orbits.
Its boundary has coefficient zero at every $\bar a$-fixed vertex:
all $p$ edges of an orbit have the same incidence sign at such
a vertex, and their total coefficient is zero in $\mathbb F_p$.

Take $h\in Q$ centralizing $\bar a$, and a word $t\in E$ lifting
$h$.  The path in $Y$ from $v$ to $hv$ obtained by projecting the
tree path from $v$ to $tv$ will be denoted by $P$.
Both endpoints are $\bar a$-fixed.  The commutator
$ata^{-1}t^{-1}\in K$ has homology class
$(\bar a-1)P$; this follows by concatenating the four projected
paths and using that $a$ fixes $v$.
A commuting lift in $R_p$ exists precisely when a cycle $z$ can
be added to $P$ so that $P+z$ is $\bar a$-invariant.
The preceding boundary observation then forces $hv=v$.
The stabilizer of $v$ in $Q$ is $q(A_i)$, and injectivity of
$q|_{A_i}$ now forces the lift of $h$ in $A_i$ to centralize $a$.
The reverse inclusion has literal lifts in $A_i$, proving the
centralizer statement.

For the normalizer replace the equation $ata^{-1}t^{-1}=1$ by
$at a^{-j}t^{-1}=1$, where $j\in\{1,\ldots,p-1\}$ is the
prescribed power in the cyclic subgroup.  Both $a$ and $a^j$ fix
$v$, and the same path calculation is $(\bar a-1)P$.
The identical argument forces $h\in q(A_i)$, where injectivity
again identifies precisely the ordinary normalizer.

Finally let $x\in R_p$ have order $p$ and nontrivial image
$h\in Q$.  Suppose that $h$ fixes no vertex of $Y$.
The subgroup $E_h=q^{-1}(\langle h\rangle)$ then has no nontrivial
vertex stabilizers in $T$: any such stabilizer would map injectively
to $\langle h\rangle$, so its generator would yield an $h$-fixed
vertex.  It also has no edge stabilizers.  Thus $E_h$ is free.
Choose the isomorphism $\langle h\rangle\simeq\mathbb Z/p$ sending
$h$ to $1$.  The resulting homomorphism $E_h\to\mathbb Z/p$ lifts
to a homomorphism $E_h\to\mathbb Z/p^2$, by choosing lifts of the
images of a free basis.  It kills $K'$ and $K^p$: elements of $K$
have image divisible by $p$.  Hence it descends through the
preimage of $\langle h\rangle$ in $R_p$.
The image of $x$ is congruent to $1$ modulo $p$ and therefore has
order $p^2$, contradicting the order of $x$.
So $h$ fixes a vertex.  Its stabilizer is a conjugate of $q(A_i)$,
and the injective restriction to $A_i$ supplies the claimed
order-$p$ element.
\end{proof}

For the three actual orbifold groups in E.256, formula (FG.1)
is, successively,
\[
\begin{array}{c|c|c}
 E&d&|R_2|\\ \hline
 C_2*C_2*C_2*F_4&1+9|Q|/2&|Q|2^{1+9|Q|/2}\\
 C_2*C_2*F_5&1+5|Q|&|Q|2^{1+5|Q|}\\
 C_2*F_6&1+11|Q|/2&|Q|2^{1+11|Q|/2}.
\end{array}                                                   \tag{FG.3}
\]
A prescribed finite quotient can be made faithful and can retain
all distinct named involution factors by adjoining their actual
abelianization map to $(\mathbb Z/2)^k$.  The resulting image has
order at most $2^k$ times the original order.  Thus (FG.3) gives
explicit bounds for the cyclic lift-class and centralizer steps
at any supplied lower quotient.

The full labelled translation subgroup $(\mathbb Z/3)^2$ in the
nine-point ambient mapping class group is a separate input.  The
lemma supplies its subsequent cyclic fixed-point steps; it does
not assert a finite conjugacy classification of that rank-two
subgroup in a finite quotient of the ambient mapping class group.
